% ============================================================================
% October 9, 2026 build (agent draft). Model section rewritten from
% 03b_model.tex in the notation of the JMP deck of Oct 9
% (latex/JMP_slides/JMP_slides.tex, frames "Environment" to "Equilibrium" and
% "Utility Specification"). Economic content is the engine of the 14.402
% base point (one-birth Estate A + net-estate floor, balanced property-tax
% rebate, PSID wealth-levels entry, post-interest transaction timing). The
% supply curve is written as in the deck, H = H0 P^eta, with H0 = 6.805: the
% engine writes H0_code (u P / r_bar)^eta with u the per-period user-cost
% rate and r_bar = 0.16, so H0 = H0_code (u / r_bar)^eta (commit feaacb00);
% the policy drivers rescale r_bar with u, so the curve in P is unchanged
% when the property tax changes.
% ============================================================================
\providecommand{\mocktodo}[1]{\textcolor{red}{[TODO: #1]}}

\section{Model}
\label{sec:model}

Motivated by these observations, I build a life-cycle model in which the same
households choose when to have children, how much housing to occupy, and
whether to rent or own it. Households live in an overlapping-generations
economy, face persistent earnings risk, and save in a single bond. Children
raise the household's needs for consumption and, above all, for space. Space
can be rented or bought. A room costs the same per period whether it is
rented or owned, but the largest dwellings are available only to owners,
owners enjoy their home somewhat more, selling a home is costly, and buying
one requires a down payment, which the household pays out of its wealth and
the income of the period. House prices and rents are determined in
equilibrium by an upward-sloping supply of housing. The government runs a pay-as-you-go pension system and taxes
property, and it returns the property-tax revenue to households as an equal
transfer.

\subsection{Environment}

Time is discrete and a period lasts four years. I describe the households,
their preferences, earnings, housing, the financial market and the
government in turn.

\subsubsection*{Demographics.}
A household enters the economy at age 18 and lives at most through age 85.
Let $a$ index its age, $a_R$ the age of retirement, 66, and $J$ the last age.
It works until $a_R$ and then receives a pension. From retirement on, it
survives from one period to the next with probability $s_a<1$; before
retirement $s_a=1$.
% [MODEL INPUT] ages 18,22,...,82 (17 cells), retirement at cell 12 (age 66),
% survival 0.9391, 0.9185, 0.8850, 0.8300 for cells 12-15, certain death after
% age 82.
\mocktodo{name the source of the survival probabilities (not documented in the
repository).}
Each household is one decision unit; I do not model the formation or
dissolution of couples.

Let $n_t$ denote the number of children ever born to the household and $m_t$
the number currently living at home. Births are possible from 18 through
age 45, at most one per period. Each child at home leaves independently with
probability $\lambda$ per period, so that a child stays on average
$1/\lambda$ periods; departure lowers $m_t$ but not $n_t$. I keep track of
$n_t$ up to three, with the last state standing for three or more. Children
are dependents: they do not work and do not make decisions.
% [MODEL INPUT] fertile cells ages 18-42 (decision at the start of each
% four-year period, so births up to age 45); lambda = 2/9 per period, mean stay
% 18 years; top state n = 3 = "3+".

\subsubsection*{Preferences.}
A household values consumption $c_t$, housing services $\hat h_t$, and
children at home. Its flow utility is
\begin{equation}
 u_t(c_t,\hat h_t;m_t)
 =\frac{1}{1-\sigma}\left(\frac{c_t^{\alpha_0}\,\hat h_t^{\,1-\alpha_0}}{e(m_t)}\right)^{1-\sigma}
 +v(m_t),\qquad v(m_t)=\psi_t\, m_t^{1-\gamma}.
 \label{eq:m-utility}
\end{equation}
The first term is utility from the consumption--housing composite per
equivalent adult. The equivalence scale
\begin{equation}
 e(m_t)=\left(\frac{2+\omega\, m_t}{2}\right)^{\nu}
 \label{eq:m-escale}
\end{equation}
measures how many adult equivalents the household contains relative to a
childless couple, so that the same spending buys less utility per person in a
larger family (as in Scholz, Seshadri and Khitatrakun, 2006); it depends on
the children at home, $m_t$, not on the children ever born. The second term,
$v(m_t)$, is the direct benefit of children at home, with $\psi_t>0$ its
level and $\gamma\in[0,1)$ the rate at which the benefit of an additional
child declines.
% [CODE] shared.py 290-312: c_bar = 0 under eqscale, escale =
% (((2+0.7 m)/2)^0.7)^(sigma-1) with m = children at home (cs), which is
% e(m)^(sigma-1), the multiplier that makes u(x/e) = e^(sigma-1) x^(1-sigma)/(1-sigma);
% child_preferences.py 31-33: benefit = psi_child m^(1-curvature) for m >= 1. The level $\psi_t$
carries a time subscript because it is the object that changes in the
historical exercise of Section~\ref{sec:quant-transition}. Becoming a parent
carries, in addition, a one-time utility cost $\xi$ at the first birth.

Children need room. A household with children at home obtains services only
from the space in excess of a minimum $\bar h(m_t)=h_P\,\1\{m_t>0\}$:
\begin{equation}
 \hat h_t=
 \begin{cases}
 h^R_t-\bar h(m_t), & \text{renter of } h^R_t \text{ rooms},\\[2pt]
 \chi\bigl[h_t-\bar h(m_t)\bigr], & \text{owner of } h_t \text{ rooms},
 \end{cases}
 \qquad \chi>1 .
 \label{eq:m-services}
\end{equation}
The parenthood space requirement $h_P$ is the space a family needs before
additional rooms add to its welfare. It is measured in rooms and does not
depend on the number of children, so it is the first child that creates the
need for space. The owner multiplier $\chi$ captures the additional services
an owner obtains from a dwelling of a given size, for example from the
ability to adapt it, and applies only to the rooms above the requirement. A
household with children in a dwelling of $h_P$ rooms or fewer obtains
essentially no housing services, so in practice parents occupy more than
$h_P$ rooms, whether they rent or own.
% [CODE] owner services chi*(h - h_P) (kernels.py 922-931: ht_c = hsv -
% owner_h_bar_scale*hbc, clamped at 1e-10 when h <= h_P, then times
% owner_service_premium; strict_hbar_feasibility is never passed, so the
% option is dominated, not infeasible). Renter: rent on h_P is paid first,
% dc = cbc + ri*hbc, and ht_unc = (1-al)*surplus/ri (kernels.py 726-729,
% 833-855). h_bar = hbar_first_child_jump + hbar_child_rooms*m (shared.py
% 291-296). The deck's "Utility Specification" frame writes
% chi[h - hbar(m)] + [h^R - hbar(m)] as one expression; read tenure by tenure,
% it is (eq:m-services).
Together, \eqref{eq:m-utility} and \eqref{eq:m-services} define what a child
costs. At given prices, a first child raises the spending needed to reach a
given utility in two ways: through the equivalence scale, which scales up all
spending, and through the space requirement, which must be paid at the rent
of a room before any housing yields services. The second part is paid every
period the child is at home, by renters and owners alike.

\subsubsection*{Fertility.}
At the start of each fertile period, a household chooses whether to try for a
child. An attempt succeeds with probability $\pi_a$, which declines with age,
and adds one child: $n_{t+1}=n_t+1$. The decision is subject to type-I
extreme-value taste shocks $\varepsilon^F_t$ with scale $\kappa_1$ when the
household is childless and $\kappa_C$ when it already has children, so that
childlessness and family size can respond differently to the same change in
incentives. After the outcome of the attempt is known, the household chooses
its housing and saving for the period. Housing can therefore adjust to a
birth within the same period, which matches the timing of the room response
in Section~\ref{sec:empirical-evidence}.
% [MODEL INPUT] pi_a = 1 - 0.02 exp(0.134 (age - 18)), zero from 45
% (parameters.py 21-45).
% [CODE] household.py 1000-1016: V = birth_count_menu(VI, n, m, pi_j,
% kappa_fert if n == 0 else kappa_fert_continuation, first_birth_fixed_cost,
% cap = 1), where VI is the housing/saving value after the birth outcome;
% birth_count.py 46-84: with cap 1 the menu is {wait, attempt} and the
% attempt value is pi*(VI[n+1,m+1] - xi*1{n=0}) + (1-pi)*VI[n,m], which is
% (eq:m-wait-attempt) and (eq:m-attempt).

\subsubsection*{Firms and earnings.}
Competitive firms produce with effective labor, $Y_t=A L_t$, where
$L_t=\int_{a<a_R} e_a z_t\,\dd G_t$ and $G_t$ is the distribution of
households over states, so the wage per efficiency unit is $w=A$, which I
normalize to one. Before retirement, a household of age $a$ earns, after the
payroll tax,
\begin{equation}
 y^d_t(a,z_t)=(1-\tau^w)\,w\,e_a\,z_t ,
 \label{eq:m-earnings}
\end{equation}
where $e_a$ is a deterministic age profile and $z_t$ an idiosyncratic labor
efficiency. Efficiency follows a first-order autoregressive process in logs,
\begin{equation}
 \log z_{t+1}=(1-\rho_z)\mu_z+\rho_z\log z_t+\epsilon_{t+1},\qquad
 \epsilon_{t+1}\sim N(0,\sigma_z^2),
 \label{eq:m-ar1}
\end{equation}
where $\mu_z$ normalizes mean efficiency to one. Earnings risk is the only
source of income uncertainty. From $a_R$ on, every retiree receives the same
pension $\varpi_t$.
% [MODEL INPUT] rho_z = 0.735, sigma_z = 0.484 per four-year period; nine
% states; pension flat.
% [CODE] e5f_social_security.py 40-41: income[:, :J_R] = (1-tax)*wage*profile,
% retirees receive the pension; shared.py 114-131: income_at_state = y*z before
% retirement, the flat pension after (retirement_income_z_scale = 0), plus the
% rebate at every age.
\mocktodo{name the data source of the age profile $e_a$ (not documented in
the files read).}

\subsubsection*{Housing and tenure.}
Housing is measured in rooms and trades at price $P_t$ per room. An owner
lives in a house of $h_t\in\mathcal H=\{2,4,6,8,10\}$ rooms. A renter rents
any $h^R_t\in[0,\bar h^R]$ rooms at rent $r_t$ per room, with
$\bar h^R<\max\mathcal H$, so the largest dwellings are reachable only
through ownership, as in the American Housing Survey facts of
Section~\ref{sec:empirical-evidence}. I call this tenure segmentation. Owners
pay depreciation $\delta_H$ and property tax $\tau^p_t$ on the value of their
home each period. Selling a home costs a fraction $\psi^s$ of its value;
buying costs nothing beyond the price. The tenure and dwelling choice is
subject to type-I extreme-value taste shocks $\varepsilon^H_t$ with scale
$\kappa_T$, one for each option: renting, keeping the current home, or buying
each size in $\mathcal H$.
% [CODE] household.py 194-206: hcost = P h, heq = (1-psi) P h, ocst =
% (delta + tau_H) P h, dp = (1-phi) P h, bmo = -phi P h; kernels.py 383-474
% (tenure_logit_kernel): options tn = 0 rent, tn >= 1 own size tn, with the
% stay value used when tn equals the current size.

\subsubsection*{Rents and housing supply.}
Competitive landlords own the rental stock. A landlord buys a room at $P_t$,
pays depreciation and property tax, rents it out at $r_t$, and sells it at
$P_{t+1}$. Landlords finance themselves at the household bond return $R_b$,
so zero profits require that rent equal the user cost of housing,
\begin{equation}
 r_t=(R_b+\delta_H+\tau^p_t)P_t-P_{t+1},
 \qquad\text{and in a steady state}\qquad
 r=(R_b-1+\delta_H+\tau^p)P .
 \label{eq:m-rent}
\end{equation}
Because rent equals the user cost and households and landlords borrow and
save at the same rate, a room costs the same per period whether it is rented
or owned. Tenure matters for the cost of a family through four features
only: the size limit on rentals $\bar h^R$, the owner multiplier $\chi$, the
cost of selling $\psi^s$, and the down payment described next. Along a
transition, \eqref{eq:m-rent} prices rent net of the expected capital gain,
so when the price is expected to recover, rent lies below the static user
cost.
% [CODE] steady state: inputs.py 159, user_cost_rate = (R_gross - 1) + delta
% + tau_H = 0.18028 per period; equilibrium.py 641, r = user_cost_rate * p.
% The dated rule along a transition is implemented in the transition runner,
% which is not in the engine files reviewed; it is stated as in the deck.

Housing is supplied by a competitive construction sector with a constant
elasticity with respect to the price,
\begin{equation}
 H^s_t=H_0\,P_t^{\eta},
 \label{eq:m-supply}
\end{equation}
where $H_0$ sets the scale of the stock and $\eta$ is the supply elasticity.
In the baseline, the stock is on this curve at every date: it adjusts to the
price within a period. Section~\ref{sec:policy-slow} replaces this with a
stock that adjusts slowly toward the same curve.
% [MODEL INPUT] eta = 0.63; H0 = 6.805 on this scale (engine: H0_code 6.312
% on the scale (u P / 0.16)^eta).
% [CODE] distribution.py 3213-3223 and 3256-3257, native_phase_b.py 97-99:
% supply = H0 * (user_cost_rate * p / r_bar) ** xi_supply, xi_supply = 0.63,
% r_bar = 0.16; demand = renter rooms + owner rooms. H0 = 6.312 *
% (0.18028/0.16)^0.63 = 6.805. inputs.py: eta_supply = 1.75 is an unused
% field; the elasticity in the law is xi_supply.

\subsubsection*{Saving and mortgages.}
Households save in a one-period bond $b_{t+1}$ with gross return $R_b$, fixed
in the world market; $b_{t+1}<0$ is mortgage debt. Borrowing must be secured
by a home: a renter cannot borrow, $b_{t+1}\geq0$. At origination, a buyer
must put down at least a fraction $1-\phi$ of the price,
\begin{equation}
 b_{t+1}\geq-\phi\,P_t h_t ,
 \label{eq:m-ltv}
\end{equation}
and afterwards the house is collateral for an owner who stays,
$b_{t+1}\geq\min\{b_t,-\phi P_t h_t\}$: a stayer pays the interest on its
mortgage and does not borrow more. An owner whose debt exceeds
$\phi P_t h_t$ after a fall in the price is therefore not forced to repay,
but cannot borrow more. The down payment is paid out of the wealth the
household brings into the period and the income it earns during it. With the
budget constraint below, \eqref{eq:m-ltv} makes a purchase feasible when
$R_b b_t+y^d_t+T_t\geq(1-\phi)P_t h_t$, so a household with little wealth
can buy out of a period's income, and there is no separate rule for how the
down payment is financed. Two further restrictions keep estates
nonnegative. From retirement on, when the household may die, its debt may not
exceed the sale value of its home net of the selling cost,
$b_{t+1}\geq-(1-\psi^s)P_t h_t$. An owner may sell and return to renting only
if the net proceeds repay the mortgage, $R_b b_t+(1-\psi^s)P_t h_{t-1}\geq0$.
Mortgages carry no amortization schedule and no limit on payments relative
to income.
% [CODE] renter floor: household.py 47-58 (unsecured_credit_limit = 0, so
% b' >= 0). Origination: household.py 205-206 (dp = (1-phi) P h, bmo =
% -phi P h), kernels.py 955-959 (total_floor = bf + unsecured_floor, with the
% unsecured part zero); phi = 0.80. Down payment from post-interest wealth
% plus income: household.py 926-931, dp_choice = (dp_arr - income)/R and
% bmo_purchase = (bmo - income)/R; kernels.py 433-445, a buyer from renting
% is infeasible if b < dp_choice or b - P h/R < bmo_purchase, both of which
% read R b + y >= (1-phi) P h (income includes the rebate, shared.py
% 114-131). Stayer: kernels.py 976-979, total_floor = max(min(b, -phi P h),
% due_death_floor); household.py 61-72. Net-estate floor: household.py
% 40-44, -(1-psi) P h when death_possible_at_age (household.py 33-38: last
% cell or survival < 1, so from the retirement cell on); kernels.py 980-983
% applied to buyers and stayers. Sale solvency: kernels.py 417-420, selling to
% rent is infeasible if R b + (1-psi) P h < 0 (require_owner_sale_solvency;
% credit.py 5-19). A seller who rebuys needs R b + (1-psi) P h_old + y >=
% (1-phi) P h_new (kernels.py 447-451). No amortization, no PTI
% (household.py 729-742 forbid them with native_purchase_income).

\subsubsection*{Budget constraint.}
Transactions in housing take place after interest accrues on the bond carried
into the period. A household that enters with $h_{t-1}$ owned rooms (zero for
a renter) and leaves with $h_t$ owned rooms and $h^R_t$ rented rooms faces
\begin{equation}
 c_t+b_{t+1}+r_t h_t^R+(\delta_H+\tau_t^p)P_t h_t
 =R_b b_t+y_t^d(a,z_t)+T_t-\1\{h_t\neq h_{t-1}\}\,P_t\bigl[h_t-(1-\psi^s)h_{t-1}\bigr],
 \label{eq:m-budget}
\end{equation}
where $y^d_t$ is earnings after the payroll tax or the pension and $T_t$ is
the lump-sum transfer from the government. A renter has $h_t=0$, an owner
$h^R_t=0$, and a household that keeps its home pays no transaction. The sale
proceeds and the purchase price are not compounded: a household that sells
and buys within a period carries only the difference into next period's
interest.
% [CODE] post-interest timing (author-adopted Oct 3): household.py 317,
% Rv = R b + y with y including the rebate; kernels.py 417 (sale to rent, b +
% (1-psi) P h/R), 433 (purchase, b - P h/R), 447 (sell and rebuy, b +
% ((1-psi) P h_old - P h_new)/R), so transactions are added to the bond after
% interest and carried into next period's interest only as a net amount.
% Owner cost (delta + tau_H) P h each period: household.py 200.

\subsubsection*{Bequests.}
A household that dies leaves its estate net of the selling cost,
$w^e_{t+1}=b_{t+1}+(1-\psi^s)P_t h_t$, and values it according to
\begin{equation}
 B(w^e_{t+1})=\theta_0\,\frac{(\theta_1+w^e_{t+1})^{1-\sigma}}{1-\sigma},
 \label{eq:m-bequest}
\end{equation}
where $\theta_0$ governs the strength of the bequest motive and $\theta_1$
the extent to which bequests are a luxury.%
\footnote{The estate enters $B$ as $\max\{w^e_{t+1},0\}$. Under the
borrowing constraints above it is nonnegative at every age at which death is
possible, so the truncation never binds.}
Estates are not received by other households; entrants' wealth is drawn from
the data instead, as described below.
% [CODE] Estate A: household.py 845-856, Vbq = bequest_utility_vec(b' + hv)
% with hv = (1-psi) P h (bequest_utility_net_active); household.py 1548-1582,
% b = max(b, 0), scale theta0*(1 + theta_n nk) with theta_n = 0, utility
% (theta1 + b)^(1-sigma)/(1-sigma), the zero-estate normalization only if
% normalize_bequest_utility (default False; the earlier footnote claiming the
% engine subtracts theta0 theta1^(1-sigma)/(1-sigma) could not be verified
% and was removed). No estate tax, estate_receiver = 'none'. Death branch:
% household.py 868-869 (last cell) and 879-881, Vnr = s*Vnr + (1-s)*Vbq.

\subsubsection*{Government.}
The government runs two programs, each with a balanced budget in every
period. The pension system pays every retiree the same pension $\varpi_t$,
financed by the payroll tax $\tau^w$ on earnings. The property tax raises
$\tau^p_t$ times the value of the entire housing stock, owned and rented:
owners pay it directly, and renters pay it through the rent, into which the
landlord passes it by \eqref{eq:m-rent}. Its revenue is returned to all
households, of every age and tenure, as an equal lump-sum transfer $T_t$, the
rebate.
% [CODE] shared.py 595-610: revenue = tau_H * price * (rental rooms + owner
% rooms); fiscal_closure.py 59-69: T*N = revenue - grants, equal lump sum at
% all ages (property_tax_rebate_closure = 'balanced'); shared.py 114-131 adds
% T to income at every age. Pension: e5f_stationary_paygo.py 47-62 solves the
% flat pension that balances the payroll tax given the exogenous age-income
% mass; e5f_social_security.py 50-96 certifies the balance (PAYGO_TOL 1e-6).

\subsubsection*{Entry.}
Households enter at age 18 as childless renters. Each entrant draws its
efficiency from the stationary distribution of $z$. Its wealth comes from the
distribution of net worth among childless renters aged 18 to 24 who head their
household in the PSID, 1984 to 2019, with negative net worth set to zero.
Dollar amounts are converted to model units with one common scale, chosen so
that mean income in this sample equals mean earnings of entrants in the
model, and observations are assigned to efficiency states by income rank, so
that the poorest entrants in the data become the least productive entrants in
the model. The children born in the model form the entering households: the
number of entrants in each period equals the number of children born sixteen
and twenty years earlier, half from each cohort, divided by 2.1, so that the
population is constant when each household has 2.1 children on average.
% [CODE] adult_entry.py 14-32: entry = adjusted births / 2.1, where a birth
% into the top state n = 3 counts tfr_top_bin_weight = 3.602 children;
% native_phase_b.py 265-270 and distribution.py 1446-1457: split queue, half
% due 16 years after birth (4 periods) and half 20 (5 periods).
% distribution.py 597-613: entrants are childless renters at the first cell,
% z drawn from z_weights, wealth from the PSID law. entry_wealth.py 197-293:
% levels_rank law, negative net worth censored at zero, one scale lambda =
% model mean annual gross entry earnings (transfers excluded) / weighted mean
% sample family income, exact weighted income-rank allocation to z states.

\subsection{The household's problem}
\label{sec:m-household}

The state of a household at the start of a period is its age $a$, bond
position $b_t$, housing $h_{t-1}$, efficiency $z_t$, and children $n_t$ and
$m_t$. A period has two stages. The household first sees its fertility taste
shocks and decides whether to attempt a birth. It then learns whether the
attempt succeeded, sees its tenure taste shocks, and chooses tenure,
dwelling, consumption and saving. Children leave home and efficiency is drawn
between periods.

\subsubsection*{Housing and saving.}
After the birth outcome, the value of the household is
\begin{multline}
 V^H_t(b_t,h_{t-1},z_t,a,n,m)
 =\E_{\varepsilon_t^H}\max_{h_t,h_t^R,b_{t+1}}
 \Bigl\{u_t\big(c_t,\hat h_t;m\big)+\varepsilon_t^H(h_t)\\
 +\beta\,\E_t\Bigl[s_a V_{t+1}(b_{t+1},h_t,z_{t+1},a+1,n,m_{t+1})
 +(1-s_a)B(w^e_{t+1})\Bigr]\Bigr\}
 \label{eq:m-vh}
\end{multline}
subject to the budget constraint \eqref{eq:m-budget}, the services
\eqref{eq:m-services}, and the borrowing constraints: $b_{t+1}\geq-\phi
P_th_t$ for a household that buys or rents, $b_{t+1}\geq\min\{b_t,-\phi
P_th_t\}$ for an owner who stays, and, when death is possible,
$b_{t+1}\geq-(1-\psi^s)P_th_t$. The expectation is over next period's
efficiency $z_{t+1}$ and the number of children $m_{t+1}$ still at home. An
option is available only if it leaves room for a saving choice that
satisfies the borrowing constraints, and a sale only if the net proceeds
repay the mortgage. Given the option $o$ (renting, keeping
the home, or buying each size), let $V^o_t$ denote the value of the best
choice of $c_t$, $h^R_t$ and $b_{t+1}$. With the extreme-value shocks,
\begin{equation}
 V^H_t=\kappa_T\log\sum_{o}\exp\bigl(V^o_t/\kappa_T\bigr),
 \qquad
 \pi^H_{t,a}(o)=\frac{\exp(V^o_t/\kappa_T)}{\sum_{o'}\exp(V^{o'}_t/\kappa_T)} .
 \label{eq:m-tenure}
\end{equation}
Because the composite in \eqref{eq:m-utility} is Cobb--Douglas, a renter not
at the size limit spends a share $1-\alpha_0$ of its outlay net of the space
requirement on rooms above $h_P$.
% [CODE] kernels.py 383-474 (tenure_logit_kernel): VH = best_v +
% kappa*log(sum exp((V^o - best_v)/kappa)) over the feasible options, which is
% (eq:m-tenure); kernels.py 833-855: renter ht_unc = (1-al)*surplus/ri, capped
% at hR_max - h_bar.

\subsubsection*{Fertility.}
Before the birth outcome, a household with $n_t<3$ in a fertile period
compares waiting and attempting, with values
\begin{equation}
\begin{aligned}
 W_t^0&=V^H_t(b_t,h_{t-1},z_t,a,n_t,m_t),\\
 W_t^A&=\pi_a\Bigl[V^H_t(b_t,h_{t-1},z_t,a,n_t+1,m_t+1)-\xi\,\1\{n_t=0\}\Bigr]
 +(1-\pi_a)\,V^H_t(b_t,h_{t-1},z_t,a,n_t,m_t).
\end{aligned}
 \label{eq:m-wait-attempt}
\end{equation}
With scale $\kappa(n_t)=\kappa_1$ if $n_t=0$ and $\kappa_C$ otherwise, the
value at the start of the period is
\begin{equation}
 V_t(b_t,h_{t-1},z_t,a,n_t,m_t)
 =\E_{\varepsilon_t^F}\max\bigl\{W_t^0+\varepsilon^F_{t,0},\;W_t^A+\varepsilon^F_{t,A}\bigr\},
 \label{eq:m-fertility}
\end{equation}
and the probability of attempting a birth is
\begin{equation}
 \pi^F_{t,a}(b_t,h_{t-1},z_t,n_t,m_t)
 =\frac{\exp\{W_t^A/\kappa(n_t)\}}{\exp\{W_t^A/\kappa(n_t)\}+\exp\{W_t^0/\kappa(n_t)\}} .
 \label{eq:m-attempt}
\end{equation}
Outside the fertile ages, or with $n_t=3$, $V_t=V^H_t$. A birth is therefore
attempted when the value of the household with one more child, net of the
cost of becoming a parent, exceeds the value of waiting by enough to overcome
the taste shocks. Anything that changes the resources a household can devote
to the space a child requires, or the risk around them, changes the gap
$W^A_t-W^0_t$ and with it the probability of a birth.

\subsection{Equilibrium}
\label{sec:m-equilibrium}

Let $G_t$ be the distribution of households over states and $N_t$ the number
of households. Given parameters, an initial population $N_0$ and an initial
distribution $G_0$, an equilibrium is a sequence
$\{V_t,g_t,G_t,N_t,P_t,r_t,T_t,\varpi_t\}_{t\geq0}$ such that, at every date:
\begin{enumerate}
\item \emph{Households.} Given $\{P_{t'},r_{t'},T_{t'},\varpi_{t'}\}_{t'\geq
t}$, $V_t$ solves the household problem of Section~\ref{sec:m-household}
with continuation $V_{t+1}$, and the policy functions $g_t$ (fertility,
tenure, dwelling, consumption and saving) attain it.
\item \emph{Rents.} The rent satisfies the landlord's zero-profit condition
\eqref{eq:m-rent}.
\item \emph{Housing market.} The rooms occupied by owners and renters equal
the supply \eqref{eq:m-supply},
\[
 N_t\int\bigl[h_t+h^R_t\bigr]\dd G_t=H_0\,P_t^{\eta}.
\]
\item \emph{Government.} The pension system balances,
$\tau^w w\int_{a<a_R}e_az_t\,\dd G_t=\varpi_t\int_{a\geq a_R}\dd G_t$, and the
rebate returns the property-tax revenue,
$T_t=\tau^p_tP_t\int\bigl[h_t+h^R_t\bigr]\dd G_t$.
\item \emph{Population.} Births are $N_t\int\pi_a\pi^F_{t,a}\,\dd G_t$.
Households survive to the next age with probability $s_a$, and the children
born form new households when they reach adulthood, as described under Entry.
\item \emph{Consistency.} $G_{t+1}$ and $N_{t+1}$ follow from $G_t$ and $N_t$
under $g_t$, the shocks, survival, the departure of children and entry.
\end{enumerate}
A stationary equilibrium is an equilibrium in which prices, the pension, the
rebate and $G_t$ are constant. In a stationary equilibrium the population is
constant only if households have 2.1 children on average, so the house price
is the one at which the population reproduces itself, and the population is
the one that clears the housing market at that price. Two features of the
equilibrium are worth stating. First, the price of housing links fertility to
housing supply: more births raise the demand for rooms by families, which
raises $P_t$ and $r_t$ and with them the cost of the space a child requires.
Second, because the housing stock is durable and the population responds to
fertility only with a lag of a generation, a change in fertility moves the
price of space for decades.%
\footnote{In the calibration I normalize the 2007 population to one and set
$H_0$ so that the market clears at the price consistent with replacement. In
every counterfactual $H_0$ is held fixed and the population adjusts.}
% [CODE] production/README.md: price root = birth-renewal residual; H0
% derived under population_one, held fixed (fixed_h0) in counterfactuals.
% equilibrium.py (model) 84-96: renewal_root = adjusted_births/(2.1*entry)
% - 1, implied H0 = demand/(u P/r_bar)^xi; native_phase_b.py 25-26, 126-127:
% RENEWAL_TOL 1e-6, population = supply/demand under fixed_h0;
% distribution.py 1101-1102 (survival), 1219-1264 (child departure via
% Pi_child and entrants maturing), 1419-1430 (property-tax revenue and rebate
% outlays on the current distribution).

\subsection{A graphical illustration}
\label{sec:m-illustration}

Figure~\ref{fig:m-illustration} summarizes how the equilibrium responds to a
fall in the preference for children. Summarize housing costs by $p$ and
fertility by $n(p;\psi)$, decreasing in $p$. Panel~(a) of the top row shows
the housing market: given the adult population $S$, the housing cost
$P(S;\psi)$ clears the market. Panel~(b) shows fertility as a function of the
housing cost. At $A$, housing clears at $(S_0,p_0)$ and fertility is at the
replacement rate $\bar n$. A fall in the preference from $\psi_0$ to $\psi_1$
first lowers fertility at the unchanged housing cost, from $A$ to $A'$
(middle row). With fewer children at home, housing demand falls and the
housing cost adjusts at the inherited population $S_0$, from $A'$ to $B$.
Over the following generations, smaller entering cohorts reduce housing
demand further (bottom row). Lower housing costs partly offset the decline in
the preference, and $C$ marks the new stationary outcome, at which fertility
is again at replacement with a smaller population and cheaper housing. The
quantitative transition of Section~\ref{sec:quant-transition} follows these
steps.

\begin{figure}[p]
\centering
\includegraphics[width=0.86\linewidth]{september_housing_fertility_stage_initial}\\[0.4em]
\includegraphics[width=0.86\linewidth]{september_housing_fertility_stage_impact}\\[0.4em]
\includegraphics[width=0.86\linewidth]{september_housing_fertility_stage_adjustment}
\caption{Fertility, housing costs and population: initial steady state,
impact, and demographic adjustment}
\label{fig:m-illustration}
\par\smallskip
\begin{minipage}{0.92\linewidth}
\footnotesize\textit{Notes:} Schematic. Panel (a) of each row: housing-market
clearing; panel (b): housing cost and fertility. Top: the initial steady
state $A$. Middle: a fall from $\psi_0$ to $\psi_1$ ($A\to A'\to B$). Bottom:
demographic adjustment to the new stationary outcome $C$.
\end{minipage}
\end{figure}
